\documentclass{beamer}
\usepackage[utf8]{inputenc}

\usetheme{Madrid}
\usecolortheme{default}


\title[MarafoNet]{MarafoNet: Marafone on the Net}
\subtitle{Architecture and implementation of a distributed, turn-based game}

\author[F. Bagattoni, J. Turchi]{Bagattoni Federico\and Turchi Jacopo}
\institute[UniBo]{Alma Mater Studiorum - Università di Bologna}
\date{\today}

\begin{document}

\frame{\titlepage}

\begin{frame}
\frametitle{Table of Contents}
\tableofcontents
\end{frame}

\section{Principles}\label{sec:principles}

\begin{frame}
\frametitle{Principles}
Which principles have we decided to respect in our project?
Talk about CAP theorem.
\end{frame}

\begin{frame}
\frametitle{Principles II}
what about the features of dependable systems?
Fault tolerance, availabiliy, reliability, redundancy, failover?
\end{frame}

\begin{frame}
\frametitle{Principles III}
Which components are introduced to do so? which nature?
\end{frame}

\section{Architecture}\label{sec:arch}

\begin{frame}
\frametitle{Architectural style}
Show off the architecture and discuss the relationship with the architecture styles shown in the course:
layered, dataspace and event-based.

the similarities are more: the system is loosely coupled and the components are totally
stateless (apart from the server that manages the persistent connections), all the communication
is done throught the dataspace, that's etcd. Seen that is divided in microservices, the
components can have their lifecycle managed separately and problems can be directly attributed
to such components.

\end{frame}

\begin{frame}
\frametitle{Architectural choices}
Justify architectural choices, statelessness and other.
\begin{block}{Remark}
Sample text
\end{block}

\begin{alertblock}{Important theorem}
Sample text in red box
\end{alertblock}

\begin{examples}
Sample text in green box. The title of the block is `Examples'.
\end{examples}

\end{frame}

\section{Storage choices}\label{sec:storage-choices}

\begin{frame}
\frametitle{Storage choice: \texttt{etcd}}
Etcd implements Raft, the fault-tolerant algorithm to distribute data and ensure consensus, so that
the game's data remains consistent even in case of server crashes.

Raft ensures that a cluster with $n=2f+1$ can survive $f$ failures.

\begin{block}{Remark}
Etcd \alert{is not protected by Byzantine Failures}: in Raft every node assumes that all the other
nodes are honest, so a malicious behavior may result in the breaking of the vote.
\end{block}
The best practice to (partially) solve this is to TODO
\end{frame}

\begin{frame}
\frametitle{Storage Considerations}
Considerations about the storage choices (is it good? is it bad?) and other possible choices that
could be made.
\end{frame}

\end{document}